\documentclass{article}
\usepackage{geometry}
\usepackage{graphicx} % Required for inserting images
\usepackage{amsmath, amsthm, amssymb}
\usepackage{parskip}
\newgeometry{vmargin={15mm}, hmargin={24mm,34mm}}
\theoremstyle{definition} 
\newtheorem{definition}{Definition}

\newtheorem{theorem}{Theorem}[section]
\newtheorem{lemma}[theorem]{Lemma}
\newtheorem{corollary}{Corollary}[theorem]


\title{MATH110BH Homework 4}
\date{February 2024}
\author{Boran Erol}

\begin{document}

\maketitle



\section{Problem 2}

\begin{lemma}
    Let $R = \mathbb{Z}[\sqrt{-5}]$. The ideal $I$ generated by $ 2 $ and $ 1 + \sqrt{-5} $ is not principal.
\end{lemma}
\begin{proof}
    Assume by contradiction that $ I $ is principal. Then, there's some $ a \in R $ that divides both $ 2 $ and $ 1 + \sqrt{-5} $.

    We use the norm argument. $ N(2) = 4 $ and $ N(1 + \sqrt{-5}) = 6 $, so $ N(a) = 2 $ so $ a = 1 $ or $ a = 2 $,
    both of which lead to a contradiction. 
\end{proof}

\newpage

\section{Problem 3}

$ 2 $ is not prime in $ \mathbb{Z}[i] $ since $ 2 = (1 + i)(1 - i) $ and $ 2 $ doesn't divide either of them.

\newpage

\section{Problem 4}

\begin{lemma}
    Let $p$ be a prime number such that $p \equiv 1 \pmod{4}$. Then, $p$ is not prime in $\mathbb{Z}[i]$.
\end{lemma}
\begin{proof}
    Since $p \equiv 1 \pmod{4}$, $\exists x \in \mathbb{Z}: x^{2} \equiv -1 \pmod{p} $.
    Notice that $ x \neq \pm 1 $ since that would imply that $ p $ is even.

    Let $ \pi: \mathbb{Z}[i] \xrightarrow{} (\mathbb{Z}/p\mathbb{Z}[i]) $.
    Recall that $ (\mathbb{Z}/p\mathbb{Z}[i]) $ is isomorphic to $ \mathbb{Z}[i] / (p) $.
    We'll denote values in the images of $ \pi $ by $ \bar{x} $.

    Notice $ \bar{x} \pm i$ or  isn't zero in $ (\mathbb{Z}/p\mathbb{Z}[i]) $.

    However, $ (\bar{x} - i)( \bar{x} + i) = 0$, so  $ \mathbb{Z}[i] / (p) $ isn't an integral domain.
    Thus, $ p $ isn't prime in $ \mathbb{Z}[i] $.
\end{proof}



\newpage

\section{Problem 6}

Recall that $ (a) \subseteq (b) $ if and only if $ b \mid a $.

Since $ \gcd(a,b) $ divides both $ a $ and $ b $, it contains both principal ideals.

If $ (c) $ contains $ (a) $ and $ (b) $, $ c $ divides both $ a $ and $ b $, so $ c $
divides $ \gcd(a,b) $. Then, the ideal generated by $ \gcd(a,b) $ is contained in $ (c) $.

\newpage

\section{Problem 7}

\begin{lemma}
    If $ R[x] $ is a UFD, $ R $ is a UFD.
\end{lemma}
\begin{proof}
    Recall from lecture that $ (R[x])^{\times} = R^{\times} $.
    
    Thus, $ R $ admits factorization by using the factorization in $ R[x] $.

    Let $ r \in R $ be irreducible in $ R[x] $. Then, $ r $ is irreducible in $ R $.
    Thus, we can use the factorization in $ R[x] $ to factor $ r $, and this is unique
    using the factorization in $ R[x] $.
\end{proof}

\section{Problem 8}

Notice that $ x^{9} + y^{9} = (x^{3} + y^{3})(x^{6} - y^{3}x^{3} + y^{6} $).

Also notice that $ x^{3} + y^{3} = (x + y)(x^{2} - xy + y^{2}) $.

Notice that $ x + y $ doesn't divide $ x^{2} - xy + y^{2} $ since setting $ y = -x $
makes this polynomial into $ 3x^{2} $ which isn't the zero polynomial. Similarly,
$ x + y $ doesn't divide $ x^{6} - y^{3}x^{3} + y^{6} $.

Then, using Eisenstein's in the ring $ \mathbb{C}[x,y][z] $ with $ p = x + y $, we conclude the proof.

Notice that $ x + y $ is prime in $ \mathbb{C}[x,y] $ since quotienting out by $ x + y $ produces a ring
isomorphic to $ \mathbb{C}[x] $ which is an integral domain.

\newpage

\section{Problem 10}

\begin{lemma}
    Let $R_{1} \subset R_{2} \subset \ldots $  be a chain of countably many subrings of a ring $R$ such that $R = \bigcup R_{i}$. Suppose that all the $R_{i}$ are UFDs and any prime element in every $R_{i}$ is prime in $R_{i+1}$. Then, R is a UFD.
\end{lemma}
\begin{proof}
    Let's first show that $ R $ admits factorization into irreducibles.
    Let $r \in R$. Then, $r \in R_{i}$ for some $i \in \mathbb{N}$. Then, $r = c_{1} \cdot \cdots \cdot c_{n}$ where $c_{i}$ is an irreducible element of $R_{i}$. Therefore, it suffices to show that $c_{i}$ is irreducible in $R$. Let $c_{i} = uv$ for some $u,v \in R$. Then, $u,v \in R_{j}$ for some $j$ by taking the maximum. Then, $u$ or $v$ is a unit in $R_{j}$. Since a unit in a smaller ring is a unit in the larger ring, $u$ or $v$ is a unit in $R$. Also notice that $c_{i}$ doesn't become a unit in $R$ since that would imply that it's a unit in some $R_{j}$, which contradicts our assumption that irreducibles stay irreducible.


    Now, suppose by contradiction that factorization is not unique. Then, there is some $x \in R$ that has two different prime factorizations, i.e. $x = a_{1}a_{2} \cdots a_{n} = b_{1}b_{2} \cdots b_{m}$. Since $a_{i}$ and $b_{j}$ are all in $R$, they are all contained in some $R_{i}$ and $R_{j}$. By taking the maximum of all of these indices, we get some $R_{i}$ such that all $a_{i}$ and $b_{i}$ are in $R_{i}$ and are prime in $R_{i}$. However, this contradicts the fact that $R_{i}$ is a UFD.
\end{proof}

\begin{lemma}
    The polynomial ring $\mathbb{Z}[x_{1},x_{2},...]$ in countably many variables is a UFD but not a Noetherian ring.
\end{lemma}
\begin{proof}
    Notice that $\mathbb{Z}[x_{1},x_{2},...] = \bigcup R_{n}$ where $R_{n} = \mathbb{Z}[x_{1},x_{2}, \ldots ,x_{n}]$.
    We have $R_{i} \subset R_{i+1}$. 
    Let $ f $ prime in $R_{i}$.
    Then, $ f $ is also prime in $ R_{i + 1} $ since any polynomial involving $ x_{i + 1} $ can't divide $ f $.
    Thus, by the lemma above, $\mathbb{Z}[x_{1},x_{2}, \ldots] $ is a UFD.

    Now, let $I_{n}$ be the ideal generated by $(x_{1}, \ldots ,x_{n})$. Clearly, $I_{n} \subset I_{n+1}$, but this increasing sequence of ideals never ends. Therefore, $\mathbb{Z}[x_{1},x_{2},...]$ is not Noetherian.
\end{proof}

\end{document}
